% Root/B/C/D full actual-file audits passed. Separate post-CP17 closure.
\subsection{Bounded geometric Brauer period}

\begin{theorem}[Proved: quadratic growth at bounded geometric Brauer period]
\label{D:bounded-geometric-Brauer-period}
Fix an integer $B\ge1$.  Let $G$ be one of the following groups,
over an algebraically closed field $k$ of characteristic $p>0$:
\begin{enumerate}
\item $G=Q(S,T)$, where $S$ is a normal proper connected surface,
      $T\subsetneq S$ is closed, and $U=S-T$;
\item $G=\pi_1(U)$, where
      $U=\operatorname{Spec}R-\{\mathfrak m\}$ for a complete
      excellent normal local $k$-algebra $(R,\mathfrak m)$ of
      dimension two with residue field $k$.
\end{enumerate}
For a continuous representation
$\rho:G\to\operatorname{PGL}_r(\mathbf F_q)$, $q=p^f$,
write $\operatorname{br}(\rho)\in H^2(U,\mathbf G_m)$
for the Brauer class of the induced algebraic projective-linear
torsor.  There is a constant $C_{G,B}$, independent of $r$ and $f$,
such that
\[
 \#\left(
  \left\{\rho:G\to\operatorname{PGL}_r(\mathbf F_q):
       \operatorname{ord}(\operatorname{br}(\rho))\le B\right\}
        /\operatorname{PGL}_r(\mathbf F_q)\right)
       \le q^{C_{G,B}r^2}.                         \tag{D.BP.1}
\]
The order of the zero class is one.  The scalar order $q-1$
is unrestricted.  More precisely, one fixed open normal
subgroup $H_B\subset G$ makes every representation counted
in (D.BP.1) geometrically Brauer-zero on its corresponding
cover of $U$.  Its choice can depend on $B$; no uniform
bound on $[G:H_B]$ as $B$ grows is asserted.
\end{theorem}

\begin{proof}
\emph{A finite set of realized geometric classes.}
Put $\operatorname{Br}'(U)=H^2(U,\mathbf G_m)_{\mathrm{tors}}$.
For each integer $m$ prime to $p$, the Kummer sequence makes
$\operatorname{Br}'(U)[m]$ a quotient of $H^2(U,\mu_m)$
\citep[Lemma~59.28.1, Tag 03PK]{StacksBrauerZeroD}.
These cohomology groups are finite.  In the complete-local
case this is the explicit etale-cohomology bound of
Theorem~\ref{D:local-surface-constant-H2}.  In the global
case it follows, with constant coefficients, from the proof
of Theorem~\ref{D:surface-tame-cross-characteristic-H2}:
proper birational Leray bounds $H^2(U,\mathbf F_\ell)$
by the cohomology of the fixed regular open on a smooth
projective resolution and a finite exceptional contribution;
the constant sheaf has zero ramification, so the cited
Betti bound applies.  A composition series for the constant
cyclic coefficient group gives the assertion for every $m$.
Tate twists are constant over $k$.

Let $\mathcal A_B$ be the set of classes $\alpha$ with
order at most $B$ that are \emph{actually realized} as
$\operatorname{br}(\rho)$ by some finite-field projective
representation of $G$, allowing every rank and every
$q=p^f$.  Such a class is killed by $q-1$, since its lifting
class comes from the central extension
\[
 1\longrightarrow\mathbf F_q^\times
 \longrightarrow\operatorname{GL}_r(\mathbf F_q)
 \longrightarrow\operatorname{PGL}_r(\mathbf F_q)
 \longrightarrow1.
\]
Consequently its order is prime to $p$, and
\[
       \mathcal A_B\subset
       \bigcup_{\substack{1\le m\le B\\(m,p)=1}}
                              \operatorname{Br}'(U)[m]
                                                        \tag{D.BP.2}
\]
is finite.  We make no assertion that every element of the
right-hand side is realized by a projective local system.

For each nonzero $\alpha\in\mathcal A_B$, choose one
witnessing representation
$\rho_\alpha:G\to\operatorname{PGL}_{r_\alpha}
(\mathbf F_{q_\alpha})$.  Define
\[
             H_B=\bigcap_{\alpha\in\mathcal A_B,
                            \,\alpha\ne0}\ker(\rho_\alpha).
                                                        \tag{D.BP.3}
\]
An empty intersection means $G$.  This is open and normal.
Let $V\to U$ be the connected finite Galois cover obtained
by composing $\pi_1(U)\to G$ with $G/H_B$.
Each witnessing projective torsor becomes trivial on $V$,
so every class of $\mathcal A_B$ pulls back to zero there.
Therefore every representation in (D.BP.1), not only the
chosen witnesses, has Brauer-zero restriction on $V$.
The fields and ranks of the finitely many witnesses may
differ; after they have been chosen the cover $V$ is fixed.

\emph{The fixed cover lies in the same geometric class.}
In case (1), normalize $S$ in $k(V)$ and call the result
$S'$.  Excellence makes $S'\to S$ finite; $S'$ is a normal
proper connected surface, and, with $T'$ the reduced inverse
image of $T$, its complement is exactly $V$.  Moreover this
cover was chosen through $G=Q(S,T)$, so it is tame at the
original boundary divisors.  Every finite quotient of the
$H_B$-action obtained by restricting a $G$-cover is tame at
the codimension-one points of $T'$: these points lie over
codimension-one points of $T$, and tameness is preserved
under the corresponding extension of discrete valuation
rings.  Thus the natural surjection
$\pi_1(V)\twoheadrightarrow H_B$ factors as
\[
               Q(S',T')\twoheadrightarrow H_B.
                                                        \tag{D.BP.4}
\]
For an arbitrary finite quotient of $H_B$, one may first
replace its open kernel by its core in $G$; this shows that
the assertion about restricted $G$-covers tests every
finite quotient of $H_B$.  The valuation comparison is the
same finite-normalization and tame-base-change argument as
Theorem~\ref{D:finite-normal-tame-descent}, using
\citep[Lemma~15.116.9(2), Tag 0EXY]{StacksDFiniteTameDescent}.
In particular no cover outside the original tame quotient
is used to choose (D.BP.3).

In case (2), the normalization $R'$ of $R$ in $k(V)$ is a
finite complete normal local domain.  Indeed finiteness
follows from excellence, and a finite domain over a
complete henselian local ring has only one local factor.
Its residue field is a finite extension of $k$, hence $k$,
and the given coefficient field remains a coefficient
field.  Its puncture is $V$.  The subgroup $H_B$ is then
$\pi_1(V)$, with the compatible base points understood.

Apply Theorem~\ref{D:Brauer-zero-projective-representation-growth}
to $(S',T')$ in case (1), and
Theorem~\ref{D:local-Brauer-zero-projective-growth} to $R'$
in case (2).  In (1), pullback along (D.BP.4) injects
conjugacy classes of $H_B$-representations into those of
$Q(S',T')$.  Thus there is one constant $C'$ for this fixed
cover such that the possible restricted conjugacy classes
in rank $r$ over $\mathbf F_q$ number at most
$q^{C'r^2}$.

\emph{The restriction fibres.}
Put $e=[G:H_B]$ and $F=\operatorname{PGL}_r(\mathbf F_q)$.
Fix one representative $\sigma:H_B\to F$ of a restricted
conjugacy class.  Every extension class can be conjugated
to have restriction exactly $\sigma$.  After choosing $e$
coset representatives for $H_B$ in $G$, such a homomorphism
is determined by its values on those representatives and
on $H_B$.  Hence there are at most $|F|^e$ extension
classes.  Since $|F|\le q^{r^2}$, the total is at most
$q^{(C'+e)r^2}$, which proves (D.BP.1).
\end{proof}

\begin{corollary}[Proved: simple products selected by bounded geometric period]
\label{D:bounded-Brauer-period-selected-simple-product}
In either case of Theorem~\ref{D:bounded-geometric-Brauer-period},
fix $B$ and $\epsilon>0$.  For each isomorphism type $L$ in
a collection of nonabelian finite simple groups, choose one
faithful embedding
\[
       \iota_L:L\hookrightarrow
                    \operatorname{PGL}_{r_L}(\mathbf F_{q_L}),
       \qquad q_L=p^{f_L},\qquad
                    |L|\ge q_L^{\epsilon r_L^2}.
\]
The quotient of $G$ selected by all distinct kernels of
epimorphisms $\eta:G\twoheadrightarrow L$ for which
$\operatorname{br}(\iota_L\eta)$ has order at most $B$
is topologically finitely generated.
\end{corollary}
\begin{proof}
Distinct kernels give distinct ambient projective conjugacy
classes.  The preceding theorem bounds their multiplicity
by $|L|^{C_{G,B}/\epsilon}$.  Apply exactly the generating-tuple,
Goursat and compactness argument of
Corollary~\ref{D:Brauer-zero-selected-simple-product}, with
$C_{G,B}$ in place of its constant $C$.
\end{proof}

\begin{remark}[Scope of the closure]
This is a finite-cover consequence of the Brauer-zero count,
not a second movement event.  The remaining obstruction in
this projective-representation route has unbounded geometric
Brauer period.  No bound uniform in $B$, no cross-characteristic
representation bound, and no control of general nonabelian
crowns is obtained here.
\end{remark}
